\subsection{Sweet--Parker消元与重联电场的定义}
由电场匹配得到$v_{\rm in}=D_\eta/\delta$，由连续性得到$v_{\rm in}=\VA\delta/L$。令两者相等：
\begin{equation}
 {D_\eta\over\delta}={\VA\delta\over L}
 \ \Rightarrow\ \delta^2={D_\eta L\over\VA}
 \ \Rightarrow\ {\delta\over L}=\left({D_\eta\over L\VA}\right)^{1/2}
 =S^{-1/2}.
\end{equation}
再代回连续性才有$v_{\rm in}/\VA=S^{-1/2}$。这说明$S^{-1/2}$来自电场、质量通量和Alfvén排气三个假设，不是从单独的欧姆定律推出的。

采用$A_z=-\psi$约定时，二维磁通差$\psi_X-\psi_O$的全导数为
\begin{equation}
 {\dd\over\dd t}(\psi_X-\psi_O)
 =\partial_t\psi|_X-\partial_t\psi|_O=E_z(X)-E_z(O).
\end{equation}
移动X/O点处$\grad\psi=0$，所以轨迹速度项消失。只在$E_z(O)$可忽略或已明确选择相应参考时，才能用单个$E_z(X)$表示磁通差的变化率。磁岛很多、参考O点跳变时，需要追踪拓扑和符号，不能只把所有局部峰值相加。

